This chapter presents reflections on the observations made during the work and, based on these, provides suggestions for practical implementation as well as proposals for future research.

\section{Discussion of Implementation and Results}

The implementation of the synthetic data generation and object detection pipeline was successful based on the results and objectives set for this work. Using the selected methods, tools, and parameters, it was possible to automate the generation of images that visually resembled the real images used for training. It is noteworthy that the chosen method, which produces synthetic images using the Stable Diffusion SDXL model, proved particularly suitable for recognizing people in indoor environments. This is likely because the base model has been trained primarily on a large collection of photographs featuring people, objects, and diverse environments. Therefore, the implementations where the base model is fine-tuned with only a few dozen images may be more difficult to apply to concepts that are very different from the original training data of the base model. Such concepts could include, for example, medical images or complex industrial components. Although numerous recent studies have investigated diffusion-based text-to-image models for these types of use cases, a detailed exploration of such applications is left for future research.

The results demonstrated that synthetic images, both alone and in combination with real images, can be effectively used to train object detection models for real-world data. This suggests that the generated synthetic images successfully captured the essential visual characteristics of the real domain. Relatively good and stable performance was achieved when a sufficient amount of synthetic data was combined with real images. When retraining object detection models multiple times while changing only the random seed, a higher mean AP indicates overall better performance, and a lower standard deviation indicates greater stability. As seen in Table~\ref{tab:ap50_seed_stats}, the s2500 model showed greater stability, the r568 model showed better overall performance, and the r568s2500 model combined both high performance and stability. This indicates that a larger dataset generally produces more stable results, while real images can provide a stronger performance boost even with fewer samples compared to synthetic images.

Several factors can explain the irregularities observed in the results, in addition to the stochasticity of deep learning model training. For example, the underperformance of the s2500 model could be attributed to overfitting, as the training loss decreased steadily while the validation loss began to increase during the last ten epochs. Conversely, the sudden high performance of the r568 model could be related to the quality of the added images, where even small improvements may have a stronger effect in a smaller dataset. This may also explain the low performance of r200, while performance and stability improved when a small number of real images were added to larger synthetic datasets. The added real images likely helped the model generalize better to the real domain, whereas r200 alone was too small to meaningfully improve the base model. The other irregularities can likely be attributed to the stochastic nature of training, where minor changes in the dataset, hyperparameters, or weight initialization can lead to significantly different outcomes. This is understandable considering that convolutional neural networks must learn to extract relevant features from images and classify them through high-dimensional fully connected layers, forming a highly complex system with many local minima in the loss landscape.

From the qualitative analysis, it can be observed that the YOLOv8n base model was already somewhat capable of detecting human figures, as its pretrained weights are likely familiar with monochromatic images. However, NIR images have their own specific noise characteristics, and in combination with difficult lighting conditions, some basic cases and especially occlusions and unusual poses proved challenging for the base model. The r568s2500 model showed a clear improvement in adapting to these conditions and was able to detect people in more challenging scenarios, even though some room for improvement remained in edge cases. It is understandable that more challenging cases still have room for improvement when training datasets are relatively small, but since the objective of this work was to test the effectiveness of the chosen methods and synthetic images, the use of much larger datasets is left to a possible full scale implementation.

\section{Practical Recommendations}

For practical implementation of the methods presented in this work, the following recommendations are made based on the observations and results obtained:

\textbf{Diversity of LoRA training dataset}: During image generation, it was sometimes necessary to use stronger prompt weights to reduce the repetition of certain visual elements originating from the training images or to make specific elements appear more prominently. Although a small number of images was sufficient for LoRA training in this work, greater diversity among them would be beneficial. For example, including different environments, people, and other varying factors in the training data would enable the generation of a wider variety of images, reducing the repetition of recurring features. Furthermore, although the generated images in this work primarily represented elevator spaces from a tilted top-view perspective, it is reasonable to assume that the model's performance could further improve if the training data included a variety of perspectives, such as close-ups, straight-on views, or entirely different environments like corridors or lobbies. In such cases, the importance of accurate and descriptive captioning becomes even more critical, as it helps distinguish these variations from one another.

\textbf{Recommended real-synthetic ratio}: 
The results of this work show that a sufficiently large mixed dataset can provide both good stability and performance, even with a relatively small number of real images. The suggested ratio of 5\%-20\% real images to 80\%-95\% synthetic images in training, as presented in Section~\ref{sec:synthetic_data_object_detection}, is supported by the results where the r200s1500 and r200s2500 models contained 11\% and 7\% real images respectively.


\textbf{Confidence threshold selection}: When the model maintains high F1-scores across a wide range of confidence thresholds, there is flexibility in selecting more suitable threshold depending on the needs of the application. For example, slightly increasing the confidence threshold can be useful if false positives are more problematic than false negatives in a given context. However, it is still recommended to select a threshold close to the one that maximizes the F1-score to achieve the best overall balance.

\textbf{Edge cases}: Although the fine-tuned models achieved good results in quantitative analysis, there were still some edge cases where detection performance was weaker, such as when people were lying on the floor. This small number of similar undetected cases could potentially be reduced by adding more images of such situations and by making people lying on the floor appear even more "blurred," as they did in the real images.
Therefore, it is recommended to include at least a few representative images of each edge case in the LoRA training dataset if possible. Alternatively, such variations could be introduced through advanced prompt engineering.

\textbf{Labeling}: It was observed across all fine-tuned models that they were vulnerable to classifying arms, hands, or partially visible limbs as separate persons. This is likely related to the labeling guidelines for occluded individuals, where the entire person was annotated even if only a small part, such as an arm, was visible or when significant overlap occurred. Even if this was not the original cause of the issue, it is recommended to carefully reconsider labeling guidelines to detect more precisely with less false detections.


\section{Future Works}

One of the most important directions for future research would be to investigate whether synthetic images could be automatically annotated during their generation phase. This would eliminate the need for manual annotation, allowing the images to be used immediately for training. For instance, the Diffusion Attentive Attribution Maps (DAAM) \cite{tang2022daaminterpretingstablediffusion} approach can produce heatmaps that visualize which parts of the generated image the diffusion model focused on with respect to a specific prompt, such as "person." Combining this approach with models like Segment Anything \cite{kirillov2023segment} could enable highly accurate automatic generation of bounding boxes for most synthetic images. Based on the author's experiments, the Segment Anything model performed well with the synthetic images used in this work, suggesting that it is a promising candidate for further investigation.

Another potential research avenue would be to identify a suitable simulator or game engine capable of modeling the desired environments, generating images from them, and then applying an image-to-image technique to make these images more realistic. This could produce data that is both precise (as in simulator-based imagery) and visually realistic, while simultaneously addressing the challenge of manual annotation. The image-to-image approach could be implemented similarly to this work within the SDXL framework by incorporating the previously mentioned ControlNet method. ControlNet would allow simulator-generated images to be transformed into the LoRA-trained visual style.

Additional research topics could include optimizing individual stages of the workflow, such as comparing different captioning methods (e.g., long descriptive captions versus concise summaries) or evaluating different upscaling approaches. Furthermore, future work could quantitatively assess the similarity between synthetic and real images using established metrics such as Fréchet Inception Distance (FID), Kernel Inception Distance (KID), or Inception Score (IS).


